\documentclass[letterpaper,oneside]{book}%

\newif\ifinstructor
%% uncomment one of the following, depending if you want to include the
%% instructor notes or not
\instructortrue
%\instructorfalse


\newif\ifnotes
%% uncomment one of the following, depending if you want to include the
%% instructor notes or not
\notestrue
%\notesfalse



\usepackage[left=1in,right=2.75in,top=1in,bottom=1in]{geometry}
\marginparwidth 1.75in

\usepackage{tabls}
\usepackage{booktabs}
\usepackage{amsmath}
\usepackage{amssymb}
\usepackage{amsthm}
\usepackage{amsfonts}
\usepackage{multicol}
\usepackage{enumitem}
\usepackage{microtype}
\usepackage{tikz}
\usetikzlibrary{positioning}

\usepackage{comment}

\usepackage{graphicx}
\usepackage{wrapfig}

\newcommand{\ds}{\displaystyle}

\let\oldmarginpar\marginpar
\renewcommand\marginpar[1]{\-\oldmarginpar{\raggedright\footnotesize #1}}
%\renewcommand\marginpar[1]{\-\oldmarginpar[\raggedleft\footnotesize #1]{\raggedright\footnotesize #1}}


%\usepackage[12hr]{datetime}
%\newdateformat{draftdate}{%
%\shortdayofweekname{\THEDAY}{\THEMONTH}{\THEYEAR}, %
%\THEDAY\ \shortmonthname[\THEMONTH] \THEYEAR}
%\draftdate
%\usepackage{eso-pic}
%\AddToShipoutPicture{\put(10,10){\small Draft: \today\ at \currenttime }}%--- version: \MakeUppercase{\svnInfoRevision}}}

% Instructor-specific material (answers, helps, etc.)
\ifinstructor
  \newcommand{\instructor}[1]{\marginpar{\textbf{Instructor: }#1}}
\else
  \newcommand{\instructor}[1]{}
\fi

\ifnotes
\renewcommand{\thefootnote}{\roman{footnote}}
\newcommand{\note}[1]{\footnote{#1}\marginpar{\fbox{\textbf{\thefootnote}}}}
\else
\newcommand{\note}[1]{}
\fi


\newcommand{\define}[2][]{\textbf{#2}\index{#1#2}}
\newcommand{\definemargin}[2][]{\marginpar{#2}\define[#1]{#2}}

% We aren't sure what we want to call "problems" (i.e., tasks,
% exercises, etc.), so we make some macros for that text.

\newcommand{\Prob}{Task}
\newcommand{\prob}{task}

\theoremstyle{plain}
\newtheorem{theorem}{Theorem}[chapter]
\newtheorem*{theorem*}{Theorem}
\newtheorem{lemma}[theorem]{Lemma}
\newtheorem*{lemma*}{Lemma}
\newtheorem{proposition}[theorem]{Proposition}
\newtheorem{corollary}[theorem]{Corollary}


\newtheoremstyle{box}%
{}{}% standard spacing before and after
{}% Body style
{}{\bfseries}{.}% Heading indent, font, and punctuation
{ }% space after heading
{\thmname{#1}\thmnumber{ #2}\thmnote{: #3}}% head spec

\newtheoremstyle{problem}%
{}{}% standard spacing before and after
{}% Body style
{}{\bfseries}{}% Heading indent, font, and punctuation
{1em}% space after heading
{\fbox{\thmname{#1}\thmnumber{ #2}\thmnote{: #3}}}% head spec

\theoremstyle{box}
\newtheorem{definition}[theorem]{Definition}
\newtheorem{dfn}[theorem]{Definition}
\newtheorem*{definition*}{Definition}
\newtheorem{observation}[theorem]{Observation}
\newtheorem{remark}[theorem]{Remark}
\newtheorem{example}[theorem]{Example}
\newtheorem{question}[theorem]{Question}
\newtheorem*{prep-problems}{Preparation Problems}

%\newtheorem{problem}[theorem]{Problem}
\theoremstyle{problem}
\newtheorem{problemnum}{\Prob}
\newenvironment{problem}[1][]{\begin{problemnum}[#1]}{\end{problemnum}\nopagebreak\hrule\bigskip}
\newenvironment{problemtodo}[1][]{\begin{problemnum}[#1]\marginpar{\large
      \textbf{TODO}}}{\end{problemnum}\nopagebreak\hrule\bigskip}


% Abbreviations
\newcommand{\ii}{\ensuremath{\vec \imath}}
\newcommand{\jj}{\ensuremath{\vec \jmath}}
\newcommand{\kk}{\ensuremath{\vec k}}
\newcommand{\vv}{\ensuremath{\mathbf{v}}}
\newcommand{\st}{\ensuremath{\,\mid\,}}
\newcommand{\colvec}[1]{\ensuremath{\begin{bmatrix}#1\end{bmatrix}}}
\DeclareMathOperator{\rank}{rank}
\DeclareMathOperator{\rref}{rref}
\DeclareMathOperator{\vspan}{span}
\DeclareMathOperator{\trace}{tr}
\DeclareMathOperator{\proj}{proj}
\DeclareMathOperator{\curl}{curl}
\newcommand{\RR}{\ensuremath{\mathbb{R}}}
% \vp is "vector prime" and corrects spacing when doing something like
% $\vec r'$ (which has the vector and prime almost touching).
% Instead, do something like $\vec r\vp$
\newcommand{\vp}{\ensuremath{^{\,\prime}}}


%The purpose of this code is to allow me to put lines in matrices so that I can create augmented matrices.
\makeatletter
\renewcommand*\env@matrix[1][*\c@MaxMatrixCols c]{%
  \hskip -\arraycolsep
  \let\@ifnextchar\new@ifnextchar
  \array{#1}}
\makeatother

\newcommand{\cl}[1]{  \begin{matrix}  #1  \end{matrix}  }
\newcommand{\mat}[1]{  \begin{bmatrix}  #1  \end{bmatrix}  }
\newcommand{\inv}{^{-1}}
\newcommand{\blank}[1]{\raisebox{0pt}[14pt]{\rule{#1}{1pt}}}


\usepackage{color}
\usepackage[breaklinks]{hyperref}
\hypersetup{
  pdfborder=0 0 2,
}

\begin{document}
\frontmatter
\title{Introduction to Linear Algebra}
\author{Jason Grout\thanks{Mathematics Faculty at Drake University,
    \url{jason.grout@drake.edu}} \and Theron
    Hitchman\thanks{Mathematics Faculty at University of Northern
      Iowa, \url{theron.hitchman@uni.edu}}}
\date{Typeset on \today\\
\vfill
\includegraphics[height=1.3cm]{by-sa}
\vfill}
\maketitle
\thispagestyle{empty}

\noindent\copyright{ 2012 Jason Grout and Theron Hitchman.  Some Rights Reserved.\\

\bigskip

\noindent Except where otherwise noted, this work is licensed under the Creative Commons
Attribution-ShareAlike 3.0 United States License. To view a copy of
this license, visit 
\begin{center}
  \url{http://creativecommons.org/licenses/by-sa/3.0/us/}
\end{center}
or send a letter to Creative Commons, 171 Second Street, Suite 300,
San Francisco, California, 94105, USA.

\bigskip

\noindent Please attribute this work to:
\smallskip

 Jason Grout, Mathematics Faculty at Drake University, Des Moines, IA,
 \url{jason.grout@drake.edu}, and

\smallskip

 Theron Hitchman, Mathematics Faculty at University of Northern Iowa,
 Cedar Falls, IA, \url{theron.hitchman@uni.edu}.

\vfill 

The source for this work is available at \url{http://github.com/jasongrout/linear-algebra-ibl}.
\vfill
}
\tableofcontents

\mainmatter

\chapter*{Preface}

Here is an interesting high-level view of the subject: a three-way
interaction between:
\begin{enumerate}
\item $\RR^n$, where our intuition lies
\item $\RR^{m\times n}$, where our computations take place
\item and abstract vector spaces, where we frame our problem and
  interpret results
\end{enumerate}


\chapter{Linear Combinations}

\section{Outline}
  Do each of these in $\RR^n$, $P_n$, $\RR^{m\times n}$, and formal
  power series rings.  Maybe also for $4\times 4$ pixel images
  (represented as matrices)?
  \begin{enumerate}
  \item Define linear combination.  See geometry in $\RR^n$
  \item Calculate inner product, length, angle, orthogonality, and
    projection in each space (except formal power series).  Examine
    these geometrically in $\RR^n$.  This introduces trace and
    transpose for the matrix space.
  \end{enumerate}

\section{Linear Combinations}

\begin{definition}\note{More weird examples of vector spaces on
    \href{http://math.stackexchange.com/questions/5233/vivid-examples-of-vector-spaces}{MathOverflow}
    (and
    \href{http://math.stackexchange.com/questions/37871/non-numerical-vector-space-examples}{here}
    and \href{http://math.stackexchange.com/questions/4694/what-are-some-alternative-definitions-of-vector-addition-and-scalar-multiplicati}{here}.}
  The word ``vector'' means different things in different contexts:
  \begin{enumerate}
  \item A vector in $\RR^n$ is a list of $n$ real numbers.
  \item A vector in $\RR^{m\times n}$ is a table of real numbers with
    $n$ rows and $m$ columns.  Such a table of real numbers is called
    an $m$ by $n$ \emph{matrix}.  Sometimes $\RR^{m\times n}$ is also
    written as $M_{m,n}$.
  \item A vector in $P_n[x]$ is a polynomial of degree less than or equal
    to $n$ with variable $x$.
  \item A vector in $\RR[[x]]$ is a formal power series (polynomial with
    infinite degree).
  \end{enumerate}
\end{definition}

\begin{problem}
  Do the following:
  \begin{enumerate}
  \item Give two examples of vectors in $\RR^3$ and two examples
    of vectors in $\RR^5$.
  \item Give two examples of vectors in $\RR^{2\times 3}$ and two
    examples of vectors in $\RR^{3\times 3}$.
  \item Give two examples of vectors in $P_2[x]$ and two examples
    of vectors in $P_3[x]$.  Is there a vector in $P_2[x]$ that is
    also a vector in $P_3[x]$?
  \item Give two examples of vectors in $\RR[[x]]$.
  \end{enumerate}
\end{problem}

\begin{definition}[Operations]
  In each of $\RR^n$, $\RR^{m\times n}$, $P_n[x]$, and $\RR[[x]]$, we
  can add two vectors by adding each part separately and we can multiply a
  vector by a scalar (a real number) by taking the number times each component.
\end{definition}

\begin{example}
  We show each operation.
  \begin{enumerate}
  \item For vectors in $\RR^2$, $(2,3)+(4,-2)=(6,1)$ and $3(2,-1)=(6,-1)$.
  \item For vectors in $\RR^{2\times 3}$, 
    \begin{align*}
      \mat{2&3&5\\1&-2&0}+\mat{4&-1&-3\\2&4&6}=\mat{6&2&2\\3&2&6}\text{,
        and}\\
      -2\mat{2&3&5\\1&-2&0}=\mat{-4&-6&-10\\-2&4&0}.
      \end{align*}
    \item For vectors in $P_3[x]$,
      $(1+x+2x^3)+(2x-x^2-x^3)=(1+3x+x^2)$ and
      $3(1+x-2x^3)=(3+3x-6x^3)$.
    \item For vectors in $\RR[[x]]$,
      \begin{align*}
        (1+x+x^2+x^3+\cdots) + (1-x+x^2-x^3+\cdots) =
        (2+2x^2+2x^4+\cdots),\\
        3(x+2x^2+3x^3+\cdots)=3x+6x^2+9x^3+\cdots
      \end{align*}
  \end{enumerate}
\end{example}

\begin{problemtodo}
  Examine the geometry of $c\vec v$, $-\vec u$, $\vec u + \vec v$, and $\vec
  u-\vec v$ in $\RR^2$.
\end{problemtodo}

\begin{definition}
  Let $c_1,c_2,\ldots,c_n$ be scalars (in other words,
  $c_1,c_2,\ldots,c_n$ are real numbers).  A \definemargin{linear
    combination} of vectors $\vec v_1,\, \vec v_2,\, \ldots,\, \vec
  v_n$ is a weighted sum $c_1\vec v_1+c_2\vec v_2+\cdots+c_n\vec v_n$.
\end{definition}

\begin{problemtodo}
  Do some linear combinations in each of the vector spaces
\end{problemtodo}

\begin{problemtodo}
  Represent a real problem as a vector in $\RR^4$ (like quantities of
  different merchandise).  Give a situation where you have a linear
  combination (like production from two suppliers).  Emphasize that
  having a vector with 4 elements is not a problem.
\end{problemtodo}

%%% Local Variables: 
%%% mode: latex
%%% TeX-master: "linear-algebra"
%%% End: 


\chapter{Vector Spaces}
\section{Outline}
Again, for $\RR^n$, $P_n$, $\RR^{m\times n}$, and formal power
series. Start emphasizing coordinate vectors more.
\begin{enumerate}
\item What is the span of these vectors?  Emphasize a vector space is
  a span of vectors.

\item Is something in the span?
  \begin{enumerate}
  \item solving systems of equations
  \item Matrix-vector multiplication as a shorthand way of writing
    linear combinations in $\RR^n$.
  \item Gaussian elimination, LU decomposition
  \item A solution is a linear combination of a vector in the null
    space and a particular solution
  \item polynomial interpolation
  \end{enumerate}

\item Is this set a span?
  \begin{enumerate}
  \item sets of solutions
  \item eigenvectors
  \item functions that vanish at a certain point
  \end{enumerate}

\item linear independence/dependence

\item basis, coordinate vectors
  \begin{enumerate}
  \item Uniqueness of coordinate vectors
  \item bases have same size
  \item building bases from linearly independent sets or spanning sets
  \end{enumerate}

\item changing definition of addition or multiplication (semi-log
  space example)
\end{enumerate}

\section{Span}

\begin{definition}
  The \definemargin{span} of a set of vectors is the set of all
  possible linear combinations of the vectors.  In symbols,
  \begin{equation*}
    \vspan(\vec v_1, \vec v_2, \ldots, \vec v_n) := \{c_1\vec
    v_1+c_2\vec v_2+\ldots c_n\vec v_n \st c_1,c_2,\ldots,c_n\in\RR\}
  \end{equation*}

  A \definemargin{vector space} is a span of vectors.  You can think
  of a vector space as a set of vectors in which we can always take a
  linear combination (it is \emph{closed} under linear combinations).

  A set of vectors \define{spans} a set if the set is the span of the
  vectors.
\end{definition}

\begin{problemtodo}
  We already looked at the geometry of spans of one, two, and three
  vectors in \Prob~\ref{prob:geometric spans in Rn}.  Look at a few
  more spans in the other spaces.
\end{problemtodo}

\begin{problemtodo}
  Examine the span of some linearly dependent vectors, and compare
  them with a corresponding basis.  Emphasize that some of the vectors
  were redundant.
\end{problemtodo}

\begin{problemtodo}
  Examine two spans that turn out to be the same set of vectors.
  Discuss why.
\end{problemtodo}

\subsection{Is this a span?}
\label{sec:this-span}

If a set is a span, then it is closed under linear combinations.

\begin{problemtodo}
  functions that vanish at a certain point.  Diagonal matrices.
  Triangular matrices.  See the webwork problem for more ideas.  See
  also the other linear algebra for some examples.  Either show in
  general (and come up with a spanning set), or exhibit a specific
  linear combination counterexample.
\end{problemtodo}

\subsection{Testing if something is in a span}

\begin{problemtodo}
  Test to see if something is in the span of a simple list of vectors,
  say in $\RR^2$.  What is the linear combination that gives the
  target vector?  Write a set of equations that you can solve to give
  you the answer.  Make the answer sufficiently complicated so that
  they can't just get it (like a fraction of one and a fraction of the
  other).
\end{problemtodo}

\begin{definition}
  If a vector $\vec u$ can be written as a linear combination of
  a list of vectors $B=[v_1,v_2,\ldots,v_n]$ as $c_1\vec v_1+c_2\vec v_2+\cdots+c_n\vec
  v_n=\vec u$, then $[\vec u]_B=(c_1,c_2,\ldots,c_n)_B$ is the \emph{coordinate
  vector} of $\vec u$ relative to $B$ (or ``with respect to
$B$'').\note{in this definition, I don't insist that $B$ is a basis
  for the set. This is a bit of a departure from normal terminology,
  but emphasizes the idea of coordinates earlier so that we can get
  comfortable with the term early on.}
\end{definition}

\begin{problemtodo}
  Write the coordinate vector from the problem above.  Explain what it
  means graphically.  Explain what the solution is graphically
  (intersection of the two lines)
\end{problemtodo}

\begin{problemtodo}
  Do another test like above, but this time involving much harder
  vectors, so you can't really just look at the solution and guess it.
  Write the coordinate vector too.  [Hint: set up a system of
  equations and solve.]  Explain what it means graphically in 3d.
\end{problemtodo}

Explain four ways of writing system of equations: equations, vector
equations, matrix equations, and augmented matrices.

\begin{problemtodo}
  Write the system in the previous problem in these 4 ways.
\end{problemtodo}

\begin{problemtodo}
  Make one of the vectors the zero vector.  Ask for two coordinate vectors
\end{problemtodo}

\begin{problemtodo}
  Make the vectors linearly dependent (but not zero).  Ask for two coordinate vectors.
\end{problemtodo}

\begin{problemtodo}
  Ask for a solution where there is no solution.
\end{problemtodo}

Emphasize that there will either be no solution, one solution, or an
infinite number of solutions. (if there are two solutions, then their
average is a solution, and any other point on the line connecting the
two solutions is also a solution).

\begin{problemtodo}
  Is the set of coordinate vectors a subspace for a homogeneous
  system?
\end{problemtodo}

\begin{problemtodo}
  Is the set of coordinate vectors a subspace for a nonhomogeneous system?
\end{problemtodo}

\subsection{Solving systems of linear equations}
\label{sec:solv-syst-line}
The sorts of systems we encountered in the last few \prob{}s are
systems of linear equations.  The linear equations we will be solving
are primarily to determine of a certain vector is in the span of some
other vectors, and if so, what the possible coefficients are.

\begin{definition}
  A system of linear equations is...  A linear equation is...
\end{definition}

\begin{problemtodo}
  Ask to solve two equivalent systems; one the rref of the other.
\end{problemtodo}

There are a series of operations we can do to simplify systems of
equations, but still have the same solution.

\begin{problemtodo}
  Graphically represent each of three operations.
\end{problemtodo}

Gaussian elimination algorithm, and practice it a few times.
Concentrate on following the specific sequence of steps and explain
why that helps you not undo work (zeroing) that you had previously done.

\begin{problemtodo}
  Solve a homogeneous system
\end{problemtodo}

\begin{problemtodo}
  Solve a corresponding nonhomogeneous system.  Point out that the
  solution is always a vector from the homogeneous solution plus a
  particular solution.
\end{problemtodo}

\begin{problemtodo}
  Polynomial interpolation application
\end{problemtodo}

\begin{problemtodo}
  LU Decomposition application?
\end{problemtodo}

\section{Linear Independence}

\begin{problemtodo}
  Repeat a problem where there are multiple solutions to solve for the
  zero vector.  Can we write one of the vectors as a linear
  combination of the others?
\end{problemtodo}

\begin{problemtodo}
  Do a problem where there is just one solution to a homogeneous
  equation.  Can we write one of the vectors as a linear combination
  of the others?  Are there multiple answers?
\end{problemtodo}

\begin{definition}
  A set of vectors is \definemargin{linearly independent} if ...
\end{definition}

\section{Bases}

\begin{problemtodo}
  Size of spanning set can't be too small.  Show that any spanning set
  of $\RR^3$ has at least three vectors, maybe.
\end{problemtodo}

\begin{problemtodo}
  Size of linearly independent set can't be too big.  Show that any
  linearly independent set in $\RR^3$ has at most 3 vectors.
\end{problemtodo}

\begin{definition}
  A \definemargin{basis} is a set that is both linearly independent
  and spans a set.  A basis spans the space, but doesn't have
  redundant vectors.  It is a ``smallest'' spanning set.
\end{definition}

\begin{problemtodo}
  All bases have the same size.
\end{problemtodo}

\begin{problemtodo}
  Find a basis for $\RR^2$, $\RR^3$, and some other spaces (maybe
  relying on previous problems?).  Write some coordinate vectors
  relative to these bases.
\end{problemtodo}

\begin{problemtodo}
  Uniqueness of coordinate vectors when you have a basis.
\end{problemtodo}

\begin{problemtodo}
  Building a basis from a spanning set
\end{problemtodo}

\begin{problemtodo}
  Building a basis from a linearly independent set
\end{problemtodo}

\section{Generalizing Operations}

List the properties (or maybe we should have them verify the
properties for a few spaces).

\begin{problemtodo}
  What if one of the properties wasn't true?
\end{problemtodo}

\begin{problemtodo}
  Examine the semi-log space.  What is the zero vector and additive inverses?
\end{problemtodo}

\begin{problemtodo}
  In the semilog space, was $0\vec v=\vec 0$?  Show this is true in
  general from the axioms.  Was $(-1)\vec v=-\vec v$?  Show it is true
  in general.  The idea of this problem is to help them see that the
  axioms are important for the operations to be intuitive and nice.
\end{problemtodo}

%%% Local Variables: 
%%% mode: latex
%%% TeX-master: "linear-algebra"
%%% End: 


\chapter{Linear Transformations}

We want mappings that preserve linear combinations.

\begin{enumerate}
\item Define linear transformation
\item Is this a linear transformation?
\item Pick bases, representation with matrices
\item kernel, image (nullspace, columnspace), (one-to-one; onto),
  rank, nullity
\item The preimage of a vector is a linear combination of a vector in
  the null space and a particular solution.  Relate this to the
  general/specific breakdown of a solution to a system of equations.
\item Transformations where dimensions of domain and codomain are the same
  \begin{enumerate}
  \item inverse transformations/matrices.  Invertible matrix theorem.
  \item change of basis matrix
  \item If the basis is orthonormal, the inverse transformation (matrix)
    is $A^T$.
  \item eigenvalues/eigenvectors---nice bases for linear
    transformations.  algebraic vs. geometric multiplicities.
  \item Determinant---area/volume change
  \end{enumerate}
\item composition of transformations/multiplication of matrices
  \begin{enumerate}
  \item an example problem might be showing that the $\rank(AB)\leq
    \min(\rank(A), \rank(B))$, see
    \href{http://www.math.ucsd.edu/~harel/Downloadable/The%20Linear%20Algebra%20Curriculum%20Study%20Group%20Recommendations%20-%20Moving%20Beyond%20Concept%20Definition.pdf}{LACSG article}, page 120, that the dimension of the image of $f\circ g \leq$ dimensions of the images of $f$ or $g$.
  \end{enumerate}
\item Symmetric matrices and spectral theorem, diagonalization
\end{enumerate}

\begin{problem}
  First Problem.
\end{problem}


%%% Local Variables: 
%%% mode: latex
%%% TeX-master: "linear-algebra"
%%% End: 


\chapter{Inner Products}

This is its own chapter, but could also be sprinkled in among the
other sections.  Advantages for dispersing this chapter is that it
slows down the course (so more time for retention) and it provides
some idea of naturally linearly independent vectors.  Of course, it
also distracts from the core focus on linear combinations and
transformations earlier, so that is an advantage for putting it all
here.

\section{Outline}
\begin{enumerate}
  \item Calculate inner product, length, angle, orthogonality, and
    projection in each space (except formal power series).  Examine
    these geometrically in $\RR^n$.  This introduces trace and
    transpose for the matrix space.
\item Examples where the inner product is not the dot product on $\RR^n$
\item Orthogonal complement; projecting onto a subspace
\item Least Squares
  \item When you have orthonormal bases, finding coordinates becomes
    much easier
  \item Making a basis orthonormal
  \item If the basis is orthonormal, the inverse transformation (matrix)
    is $A^T$.
\item Symmetric matrices and spectral theorem, diagonalization

\end{enumerate}

%%% Local Variables: 
%%% mode: latex
%%% TeX-master: "linear-algebra"
%%% End: 


\chapter{Projects}

This unit provides some ideas for final projects.

\begin{enumerate}
\item DCT
\item Vector fields/ODE connection (Wronskian?  Solutions to
  differential equations?)
\item Page Rank/Markov Processes
\item Software for linear algebra (Matlab, LINPACK, BLAS, MMX, cache
  stages in a CPU and block matrix arithmetic, GPU linear algebra)
\item Pseudoinverses
\item Jordan form
\item QR decomposition
\item SVD
\item Legendre polynomials as an inner product space?
\end{enumerate}

%%% Local Variables: 
%%% mode: latex
%%% TeX-master: "linear-algebra"
%%% End: 


\end{document}

%%% Local Variables:
%%% mode: latex
%%% TeX-master: t
%%% End:
